% ================================================================
\section{The Cathedral and Society}
\label{sec:society}
% ================================================================

\subsection{Open Science and the Crisis of Reproducibility}

The sciences are in the midst of a reproducibility crisis: across
psychology, medicine, and economics, a significant fraction of
published results fail to replicate when independent researchers
attempt to verify them. Mathematics has traditionally been
considered immune to this crisis, on the grounds that mathematical
proofs are ``certain'' in a way that empirical results are not.

This confidence is misplaced. Published mathematical proofs
regularly contain errors. Some are benign (notational slips that
an expert can correct); others are serious (logical gaps that
invalidate the argument). The history of the classification of
finite simple groups, Mochizuki's claimed proof of the ABC
conjecture, and the numerous retractions catalogued by
Retraction Watch demonstrate that mathematical peer review is
fallible.

The Cathedral offers a different model: \emph{formal
reproducibility}. The entire proof is available as a public
repository. Any person with a Lean~4 installation---which is
freely downloadable---can compile the repository and verify
every theorem. The verification is:

\begin{itemize}
  \item \textbf{Complete}: every line is checked, not just the
    lines a referee chooses to examine.
  \item \textbf{Deterministic}: the same source code produces
    the same verification result on any computer.
  \item \textbf{Incorruptible}: the compiler cannot be persuaded
    by reputation, social pressure, or incomplete reasoning.
  \item \textbf{Permanent}: the verification can be repeated at
    any future date, as long as the Lean compiler exists.
\end{itemize}

\noindent This is a stronger form of reproducibility than any
science currently achieves. It suggests that formalization is not
merely a niche activity for foundations enthusiasts but a model
for how all mathematical knowledge could be organized and verified.

\subsection{The Democratization of Verification}

Traditional mathematics has a high barrier to entry for
verification. To check a claimed proof of a number-theoretic
result, one needs:

\begin{enumerate}
  \item Graduate-level training in analytic number theory.
  \item Familiarity with the specific techniques used (contour
    integration, mean value theorems, character sums, etc.).
  \item Sufficient time (weeks to months) to work through the
    details.
  \item Access to the relevant journals and books.
\end{enumerate}

The Cathedral lowers this barrier dramatically. To verify the
Cathedral's claims, one needs:

\begin{enumerate}
  \item A computer.
  \item An internet connection (to download Lean and Mathlib).
  \item The command: \texttt{lake build}.
\end{enumerate}

\noindent No mathematical training is required. The compiler
does not care whether the user understands the proof; it only
requires that the proof type-checks.

This is a genuine democratization. A teenager in a rural
school, a retired engineer, a philosopher with no mathematical
background---any of these can verify the Cathedral's claims
with absolute certainty, simply by running a compiler. The
epistemic authority shifts from \emph{people} (``trust me, I'm a
Fields Medalist'') to \emph{code} (``trust the compiler, here's
the repository'').

We do not claim that this shift is unambiguously good. There
is genuine value in human mathematical understanding, and a
culture that values only machine-checkable proofs risks losing
the intuition, creativity, and aesthetic sensitivity that drive
mathematical progress. But the shift is happening, and the
Cathedral demonstrates its feasibility for problems of genuine
mathematical significance.

\subsection{The Long View: Mathematical Permanence}

Euclid's \emph{Elements} has survived 2{,}300 years. It survives
because its arguments are rigorous by the standards of each
era that encounters them, and because the results it proves
are \emph{true}---verified by each generation's own methods.

What is the permanence of a Lean proof? The answer depends on
the permanence of the Lean type theory. The Calculus of Inductive
Constructions is a mathematical object, not a software product.
Its rules are precise and can be implemented in any programming
language. Even if the Lean project ceases to exist, the type
theory persists, and a new implementation could verify the
Cathedral's proofs.

In this sense, a formal proof has \emph{greater} permanence than
an informal one. An informal proof depends on the continued
existence of a community that can read and evaluate it.
A formal proof depends only on the continued existence of
a type theory---a mathematical object that, like all mathematical
objects, exists independently of its representations.

The Cathedral may outlast not only its creators but the
civilization that built it. If a future society rediscovers the
CIC type theory and implements a type-checker, they can verify
every theorem in the Cathedral without any knowledge of its
authors, its historical context, or the Riemann Hypothesis itself.
The proof is self-contained, self-verifying, and permanent.

\subsection{What If RH Is Proved Tomorrow?}

If someone---human, AI, or collaboration---formalizes the Crown
Axioms in Lean, the Cathedral converts from a conditional result
to an unconditional one. The entire 120{,}000-line infrastructure
becomes a machine-verified proof of the Riemann Hypothesis.

This conversion is not hypothetical but \emph{designed}. The
Cathedral's architecture includes a single entry point
(\lean{nyman\_beurling\_equivalence}) that accepts the Crown Axiom
as input and produces $\RH \iff d_N^2 \to 0$ as output. Plug in
the axiom; get the Riemann Hypothesis.

The philosophical significance is that the \emph{work has been done
in advance}. The Cathedral is a mathematical structure waiting
for two inputs. When those inputs arrive---whether in weeks, years,
or decades---the proof is complete. The builders have laid the
foundation, erected the walls, installed the stained glass,
and calibrated the observatory. All that remains is to open the
door.

This ``just-in-time'' architecture is a new mode of mathematical
work: not proving a theorem, but \emph{preparing the infrastructure
for a theorem's arrival}. It reflects a confidence that the
theorem will eventually be proved---not because we know it is true,
but because we have built a vessel capable of containing it.
